\newpage
\phantomsection\addcontentsline{toc}{section}{Recursos computacionales}
\section*{Recursos computacionales}

Para esta tarea, nos apoyamos del Laboratorio de Supercómputo de CIMAT. El servidor de este Laboratorio
cuenta con dos unidades NVIDIA TITAN RTX, con 24 GB de VRAM por GPU. Además, cuenta con 12.3 nodos y 12.1
con conda. El ecceso es remoto, vía SSH y trabajamos con nodos computacionales sin acceso a internet. La 
descarga de los modleos Transformer se realizaron primero a mi computadora y luego se suvieron al Laboratorio. 

Esto nos permitió manejar y finetunear modelos más grandes de lo que mi computadora puede soportar, como 
lo fueron Mistral-7B o LLaMA. En mi caso, los modelos corridos en mi PC fueron las arquitecturas de redes 
neuronales recurrentes (RNN, LSTM y GRU). Mi computadora cuenta con una NVIDIA GeForce GTX 1650, de 4GB de 
de VRAM y 12.4 nodos. 

Entre las cosas que aprendimos, aparte de lo solicitado en la tarea, fue el manejo de SLURM y de comandos
para servidores como el de CIMAT. 

Los tiempos de ejecución para los modelos de Transformers en general fueron rápidos, aproximadamente 10 por prueba. En 
el caso de los modelos de redes neuronales recurrentes, los tiempos de ejecución en mi computadora tardaron máximo 
30 minutos por modelo, dependiendo de los hiperparámetros. 